% !TEX root = ../main.tex

\begin{abstract}[zh]
时空转录组学（Spatial Transcriptomics）的快速发展，为在单细胞分辨率下研究多细胞生物发育的动力学过程提供了前所未有的契机。然而，如何在连续时间序列的时空转录组数据中准确重建细胞的谱系分化和迁移轨迹，仍然是当前计算生物学的重大挑战。传统的谱系追踪算法（如 WOT 和 Moscot）通常基于欧氏空间的最优传输（Optimal Transport, OT）理论，这导致其在处理具有复杂空间几何拓扑（如发育空腔、物理屏障）的组织时，容易产生穿越物理真空区域的“欧氏空间瞬间移动（Euclidean Teleportation）”伪影。同时，传统的无向最优传输模型缺乏对生物时间方向的物理约束，导致细胞在复杂分叉处的动力学模拟失真。

为了解决上述关键痛点，本文提出了一个物理和几何双重约束的时空单细胞谱系追踪模型——SpaLineage-OT。该模型的核心贡献包含以下三个方面：第一，引入黎曼测地线几何约束，构建基于 Delaunay 三角剖分的空间流形图，利用核密度估计（KDE）定义空间形态阻抗函数，并通过 Dijkstra 最短路径算法计算测地线距离，以此取代传统的欧氏距离，强制细胞最优传输耦合贴合真实的组织流形，消除了跨越物理真空的“瞬间移动”伪影；第二，提出非对称速度引导传输代价（Asymmetric Velocity-Guided Cost），将单细胞 RNA 速度（RNA Velocity）的余弦相似度作为非对称方向性惩罚项融入融合格罗莫夫-瓦瑟斯坦（FGW）框架中，建立了具备时间方向箭头的最优传输模型；第三，基于连续时间生成式学习，采用薛定谔桥流匹配（Schrödinger Bridge Flow Matching, SBFM）结合 DriftMLP 神经网络拟合细胞发育的连续漂移速度场，并在推理端设计了均值漂移势能引导（Mean Shift Potential Guidance）机制与四阶龙格-库塔（RK4）常微分方程求解器，对边缘和低密度真空区域的细胞漂移向量进行物理纠偏，保证细胞在演化中不会发生数值下溢发散。

我们在小鼠胚胎器官发育时空转录组 Stereo-seq 数据集（MOSTA, 包含 E9.5、E10.5 和 E11.5 三个关键时间点）上对 SpaLineage-OT 进行了严谨的留出验证（Hold-out Validation）。定量评估结果表明，相较于基线方法 Moscot 和 WOT，SpaLineage-OT 重建的 E10.5 细胞空间密度分布与真实密度的能量距离（Energy Distance）降低了 15.5\%。在物理屏障穿越率（Barrier Crossing Rate, BCR）方面，SpaLineage-OT 达到了 91.40\%，显著优于 WOT 的 33.20\% 和 Moscot 的 45.80\%，有效防止了细胞轨迹跨越脑室等物理空腔。此外，细胞迁移向量与单细胞 RNA 速度的方向一致性也得到了显著提升。本研究表明，将生物物理先验与黎曼几何融入最优传输及连续流匹配理论，对于准确捕捉和解释复杂组织发育中的单细胞动态演化轨迹具有重要的科学意义和学术应用价值。
\end{abstract}

\begin{abstract}[en]
The rapid development of spatial transcriptomics has offered unprecedented opportunities to study the dynamic processes of multicellular organisms at single-cell resolution. However, reconstructing continuous cell differentiation lineages and spatial migration trajectories from time-series spatial transcriptomics remains a fundamental challenge in computational biology. Conventional lineage tracing algorithms (such as WOT and Moscot) are formulated based on Optimal Transport (OT) under Euclidean assumptions. Consequently, when analyzing tissues with complex spatial topologies (e.g., developmental cavities or physical barriers), these methods suffer from "Euclidean Teleportation" artifacts, where cell trajectories erroneously cross physically inaccessible vacuum regions. Furthermore, standard undirected OT formulations lack biophysical directional constraints, resulting in inaccurate trajectory predictions at complex developmental branching points.

To address these limitations, we propose SpaLineage-OT, a physically and geometrically constrained spatiotemporal single-cell lineage tracing framework. The core contributions of this work are threefold: First, we introduce a Riemannian geodesic geometric constraint. By constructing a spatial manifold graph using Delaunay triangulation, measuring morphological impedance via Kernel Density Estimation (KDE), and computing shortest paths using Dijkstra's algorithm, we replace the Euclidean distance with geodesic distance. This forces the Fused Gromov-Wasserstein (FGW) coupling to adhere to the physical tissue manifold and eliminates the "teleportation" artifacts across spatial voids. Second, we design an Asymmetric Velocity-Guided Cost, incorporating RNA velocity cosine similarity as an asymmetric directional penalty into the FGW framework, thereby establishing a directed optimal transport mapping. Third, within a continuous-time generative learning paradigm, we implement Schrödinger Bridge Flow Matching (SBFM) parameterized by a neural network (DriftMLP) to fit continuous developmental velocity fields. At the inference phase, we introduce a Mean Shift Potential Guidance mechanism coupled with a fourth-order Runge-Kutta (RK4) ODE solver to apply physical corrections to cell drift vectors in low-density or vacuum regions, preventing numerical underflow and path divergence.

We rigorously evaluated SpaLineage-OT using a hold-out validation setup on the spatiotemporal Stereo-seq dataset of mouse organogenesis (MOSTA, spanning E9.5, E10.5, and E11.5). Quantitative results demonstrate that SpaLineage-OT reduces the Energy Distance (ED) between the predicted and ground-truth spatial density distributions of E10.5 cells by 15.5\% compared to baseline methods. In terms of tissue barrier crossing rate (BCR), SpaLineage-OT achieves 91.40\%, significantly outperforming WOT (33.20\%) and Moscot (45.80\%), effectively preventing cell trajectories from artificially penetrating developmental cavities (such as brain ventricles). Additionally, the alignment between cellular migration vectors and RNA velocity directions shows substantial improvement. This study demonstrates that incorporating biophysical priors and Riemannian geometry into optimal transport and continuous-time flow matching is critical for accurately tracking and understanding single-cell dynamics during complex tissue development.
\end{abstract}
